
\subsubsection{6.1 Interpretation of Main Findings}

\textbf{Finding 1: Architectural invariance produces machine-precision zero OrbitVar.} DeepSets achieves OrbitVar = 1.002e-14, not merely small but at the precision floor of float32 arithmetic. This is not a quantitative improvement over CISE --- it is a qualitative regime change. Once OrbitVar is at machine precision, there is no residual symmetry violation for downstream predictors to compensate for; the encoder output is definitionally the same for all functionally equivalent models.

NFN's OrbitVar of 8.905e-08 represents near-invariance but not exact invariance. The small residual is consistent with NFN's HNPPool aggregation introducing finite-precision imprecision from spatial interaction terms in the CNN permutation group. Whether this residual affects downstream prediction performance requires an experiment that was blocked by NFN library installation failure (Section 4.4).

\textbf{Finding 2: CISE's non-invariance is a signal problem, not a noise problem.} MSE\_perm/MSE\_total = 3.3452 demonstrates that CISE's permutation sensitivity does not add small noise on top of a clean signal --- it contributes an error component larger than the total prediction error budget. The $R^{2}(C1_avg) = -1.6288$ result confirms the mechanism: CISE encodes channel-position information as a primary predictive feature, and averaging over functionally equivalent permuted embeddings destroys the predictor entirely by presenting inputs the predictor was never trained to handle.

\textbf{Finding 3: The entanglement result is an empirical finding rather than a failure.} The additive closure criterion $\Delta MSE \approx MSE_perm^{C1} \pm$ 10\% was pre-registered as the quantitative test of the causal mechanism. It fails decisively (closure = 0.872). This is presented as an honest result: the direction of the causal mechanism (architectural invariance $\rightarrow$ lower MSE $\rightarrow$ higher $R^{2})$ is confirmed, but the magnitude is not predictable from $MSE_perm^{C1}$ alone. The closure failure indicates that MSE\_perm and MSE\_res in CISE embedding space are not orthogonal --- they are correlated through the downstream predictor's non-linear feature interactions.

The most plausible mechanism is that LightGBM trained on CISE embeddings implicitly learns to suppress the most permutation-sensitive embedding dimensions in favor of more stable ones. This implicit compensation reduces effective MSE\_perm below its naive estimate, while simultaneously changing MSE\_res. When C2 removes channel-position encoding architecturally, LightGBM re-optimizes over a geometrically different embedding space; both MSE\_perm and MSE\_res change, and their combined change does not equal the naive prediction. This interpretation predicts that simpler predictors with less capacity for implicit suppression would show less entanglement, a hypothesis that was not tested in the present work (Ridge regression on C2 embeddings yields $R^2 = -6.42$, ruling out linear heads as a valid comparison point).

\subsubsection{6.2 Limitations}

\textbf{Additive closure not achieved (closure = 0.872).} The pre-registered quantitative prediction fails empirically. The direction of effect is confirmed; the magnitude is not. Future work should establish whether additive closure can be recovered under specific conditions (orthogonalized embedding spaces, different predictor architectures) or whether entanglement is a general feature of non-linear predictors on non-invariant encodings.

\textbf{NFN downstream comparison unavailable.} NFN achieves OrbitVar = 8.905e-08, confirming near-invariance at the representational level, but downstream $R^2$ comparison against DeepSets was blocked by library installation failure. Kendall's $\tau$ = 0.934 reported by Zhou et al. (2023) for NFN on CIFAR-10-GS generalization prediction uses a different evaluation protocol and cannot be directly compared to the $R^2$ results in this work. This comparison is the primary open empirical question.

\textbf{Single dataset and architecture.} All experiments use ModelZooDataset CIFAR10-GS: 100 CNNs with a fixed 3-layer structure (permutation group $S_8 \times S_6 \times S_{4})$. Generalization to wider networks (ResNets, ViTs), larger model zoos, or different prediction tasks requires separate experiments. The MSE decomposition diagnostic is theoretically architecture-agnostic, but specific magnitudes of OrbitVar and MSE\_perm will depend on the dataset and architecture.

\textbf{$R^2$ reference point sensitivity.} Unterthiner et al. (2020) report $R^{2}(C0)$ = 0.984 from 5-fold training-set cross-validation on the full dataset; the testset evaluation here yields $R^{2}(C0)$ = 0.7316. Matched evaluation protocol would enable direct comparison with the published baseline. The MSE\_perm mechanism (ratio = 3.3452) and the direction of the C1 $\rightarrow$ C2 improvement are evaluation-protocol-independent; Kendall's $\tau$ improvement (0.7205 $\rightarrow$ 0.7651) is also robust to protocol differences.

\textbf{Distribution-shift robustness not evaluated.} A planned experiment evaluating encoder performance on CIFAR-10-C (corrupted inputs) was not executed. Distribution-shift robustness claims for invariant encoders are not supported by the present experiments.

\subsubsection{6.3 Broader Impact}

This work establishes that weight-space encoders should be evaluated not only on downstream performance but on permutation sensitivity (OrbitVar) and its propagation to prediction error (MSE\_perm). The MSE bias-variance decomposition over permutation orbits introduced here is a reusable diagnostic that can be computed for any encoder-predictor pair with access to functional permutations.

For practitioners building model zoo applications, the result is that architectural invariance (DeepSets doubly-invariant sum pooling) yields a 6.37 percentage-point $R^2$ improvement over a non-invariant encoder at no additional training cost. No negative societal impacts are anticipated from this work.

